\chapter{Introduction to SteamCAD2}\label{chap:intro}
\special{pdf: out 1 << /Title (\thechapter. Introduction to SteamCAD2) /Dest [ @thispage /FitH @ypos ] >>}

SteamCAD2 is a 2D drawing tool especially designed to create drawings of
steam locomotives as you can see on the next page\footnote{This drawing presents the
Austrian steam locomotive class 6 kkStB. This type was build in several Austrian
factories from 1894 to 1898. The image is based on the factory's drawing and is at
exact scale 1:87 (known as H0 scale).}. First of all, why is it called SteamCAD2?
There are basically two reasons for this name. One was already mentioned - the tool is
specially designed for drawing steam engines. But there is another meaning, it should
emphasize that the software is very lightweight. The user interface is minimalistic,
the application window has a menu bar and a status bar only. There are no toolbars with
thousands of buttons and functions. The main aim when designing SteamCAD2 was that the
application contains as few tools, as few controls, as possible. All the functionality
is accessible from the menu, and each command has a shotcut, so that it is possible to
work with the application just using mouse and keyboard. For experienced user, almost
no interaction with the application GUI is required.

As a result, the drawing area is not obscured by controls and can occupy almost the whole
screen. This is what makes the SteamCAD2 different from any other CAD. SteamCAD2 is focused
on the productivity, not on the fancy look.

\section{History}
\special{pdf: out 2 << /Title (\thesection. History) /Dest [ @thispage /FitH @ypos ] >>}

SteamCAD2 is an evolution of SteamCAD. It has several enhancements, namely paths and areas
can be defined and filled, and also colors has been added.

\pagestyle{empty}
\begin{center}
\vspace*{-2.2cm}
\hspace*{-0.7cm}
\makebox[\linewidth]
{
\includegraphics[angle=90,origin=c]{Images/264_0.eps}
}
\end{center}
\newpage

\pagestyle{headings}

\section{What is SteamCAD2}
\special{pdf: out 2 << /Title (\thesection. What is SteamCAD2) /Dest [ @thispage /FitH @ypos ] >>}

SteamCAD2 is an open source software for creating 2D technical drawings, especially designed to 
% \makebox[\linewidth][s]{produce vector drawings suitable for further publishing process. The
% software itself does not}
produce vector drawings suitable for further publishing process. The software itself does not
provide any printing capabilities. Its aim is solely to produce drawings in a format consumed
by other technologies, either desktop or web based.

The software features several advanced technologies, namely length and angular unit handling, both
for entering and for printing, and linear pattern conflict resolution. Both will be explained
later in more details.

The software is available in Linux and Windows version. Both versions use native Windows
management, and different technology for drawing backend. It uses Cairo on Linux and GDI+ on
Windows.

The philosophy of SteamCAD2 is to follow the manual drawing and adding some advanced features
not available for hand drawing, such as mouse snapping tool.

\section{What is new in SteamCAD2}
\special{pdf: out 2 << /Title (\thesection. What is new in SteamCAD2) /Dest [ @thispage /FitH @ypos ] >>}

As was mentioned, SteamCAD2 has several enhancements. First of all, the drawing backend on Windows
has been changed to GDI+, thus the Windows user experience is now almost the same like on Linux.
Furthermore the line dash patern drawing was changed so that it uses the native dash capabilities
of either Cairo or GDI+.

Then an option to create path from adjacent line segments was added, which makes the rendering
of dash lines better. Finally it is also possible to define filled areas, so that more artistic
drawing can be achieved now, yet still retaining the CAD capabiities.

\section{Installing SteamCAD2}
\special{pdf: out 2 << /Title (\thesection. Installing SteamCAD2) /Dest [ @thispage /FitH @ypos ] >>}

On a GNU Build system it is possible to install SteamCAD2 in a standard way:
\begin{verbatim}
$ sudo apt-get install gcc libcairo2-dev gtk+-dev
$ git clone https://github.com/oskardolch/SteamCAD2.git
$ cd SteamCAD2
$ ./configure
$ make
$ sudo make install
\end{verbatim}

SteamCAD2 then appears in the Applications > Graphics menu.

On a Windows system, the best is to download the zip with binaries for either 32bit or for 64bit
platform. Note that Windows XP and older systems are not supported by the provided build. If
you want to run it on an unsuported Windows system, you should compile it yourself.

There is no installation utility provided for MS Windows, just unzip the package into a directory
where you have write access to. So the best choise is somewhere under your documents folder.

\section{Building SteamCAD2}
\special{pdf: out 2 << /Title (\thesection. Building SteamCAD2) /Dest [ @thispage /FitH @ypos ] >>}

For a debuging purpose it is also possible to build SteamCAD2 bypassing the GNU build system.
You should install the follwoing dependencies first: gcc, cairo-dev and gtk+-dev.
Then download the source code and compile it from the console. Before you run
the application, copy the files DPapers.ini and DUnits.ini from the Papers folder into the
/home/<me>/.SteamCAD2 folder. The meaning of the ini files will be explained later.

The compilation steps and SteamCAD2 invocation on Debian based systems might look as the following:
\begin{verbatim}
$ sudo apt-get install gcc libcairo2-dev gtk+-dev
$ git clone https://github.com/oskardolch/SteamCAD2.git
$ cd SteamCAD2
$ mkdir Build
$ make -fmakefile.gcc libdxflib.a
$ make -fmakefile.gcc SteamCAD2
$ ./Build/SteamCAD2
\end{verbatim}

Building of the software on MS Windows is not straightforward, since it involves building of
the cairo.dll.

\section{Prerequisities}
\special{pdf: out 2 << /Title (\thesection. Prerequisities) /Dest [ @thispage /FitH @ypos ] >>}

SteamCAD2 has minimal hardware requirements, you should be able to run it on almost every
computer. However, what you desperately need is a 3-button mouse with wheel. There is no way
to work with SteamCAD2 without such a device.

